% Generated. Do not edit.
% Projection scriptural-date-location/evidence, version 1: the Scriptural Date and
% Location page, generated whole from research/chronology.toml, the Scripture
% chronology corpus and the source library by tools/tpt document-types refresh.
% document: liturgy/roman-rite/1962/propers/temporal/44-fourth-after-pentecost
\section*{Appendix: Evidence for the Scriptural Dates}
\label{sec:date-location-evidence}
{\small For each passage, as the Scriptural Date and Location table names it: what is dated, how the corpus weighs each answer, the source's own words, and the sources cited. A subject the corpus does not date is reported as such and given no date.\par}
{\small
\par\smallskip\noindent\textit{Alleluia: Ps. 9:5, 10 (Vg./Heb.)}\par\nobreak
\begin{itemize}[leftmargin=1.2em,itemsep=1pt,topsep=1pt,parsep=0pt]
\item \textbf{Historical setting}. No date: the chronology corpus has not yet dated this (research pending).
\item \textbf{Composition} (preferred). Psalms: Before c.~165 B.C. (``before the Maccabean period, around 165 B.C.; no individual psalm can be dated securely''; Catholic critical chronology; New American Bible, Revised Edition --- USCCB web state (2026), Psalms introduction).
\end{itemize}
\par\smallskip\noindent\textit{Offertory: Ps. 12:4--5 (Vg.; Heb. 13)}\par\nobreak
\begin{itemize}[leftmargin=1.2em,itemsep=1pt,topsep=1pt,parsep=0pt]
\item \textbf{Historical setting}. As under Alleluia: Ps. 9:5, 10 (Vg./Heb.).
\item \textbf{Composition}. As under Alleluia: Ps. 9:5, 10 (Vg./Heb.).
\end{itemize}
\par\smallskip\noindent\textit{Communion: Ps. 17:3 (Vg.; Heb. 18)}\par\nobreak
\begin{itemize}[leftmargin=1.2em,itemsep=1pt,topsep=1pt,parsep=0pt]
\item \textbf{Superscription setting} (preferred). David's deliverance from all his enemies and from the hand of Saul: In the day that the Lord delivered him out of the hand of all his enemies, and out of the hand of Saul (The Holy Bible, Douay-Rheims (Challoner revision), Project Gutenberg eBook 1581: 2 Kings 22:1, Ps. 17:1).
\item \textbf{Historical setting} (preferred). The reign of Saul: ``Saul was a child of one year when he began to reign, and he reigned two years over Israel'' (duration) (The Holy Bible, Douay-Rheims (Challoner revision), Project Gutenberg eBook 1581: 1 Kings 13:1).
\item \textbf{Utterance} (preferred). David's deliverance from all his enemies and from the hand of Saul: In the day that the Lord delivered him out of the hand of all his enemies, and out of the hand of Saul (The Holy Bible, Douay-Rheims (Challoner revision), Project Gutenberg eBook 1581: 2 Kings 22:1, Ps. 17:1).
\item \textbf{Composition}. As under Alleluia: Ps. 9:5, 10 (Vg./Heb.).
\item \textbf{Prophetic referent} (disputed). The Nativity of Our Lord: B.C.~3 (``the year of Rome 750 which he styles 3 B.C.''; The Catholic Encyclopedia, Volume 3 (New York, 1908)); c.~B.C.~7 (``probably the year 7 B.C.''; The Catholic Encyclopedia, Volume 3 (New York, 1908)). Place: Bethlehem as the Christian referent (The Holy Bible, Douay-Rheims (Challoner revision), Project Gutenberg eBook 1581: Matt. 2:1).
\end{itemize}
\par\smallskip\noindent\textit{Introit: Ps. 26:1--2 (Vg.; Heb. 27), Ps. 26:3 (Vg.; Heb. 27)}\par\nobreak
\begin{itemize}[leftmargin=1.2em,itemsep=1pt,topsep=1pt,parsep=0pt]
\item \textbf{Historical setting}. As under Alleluia: Ps. 9:5, 10 (Vg./Heb.).
\item \textbf{Composition}. As under Alleluia: Ps. 9:5, 10 (Vg./Heb.).
\end{itemize}
\par\smallskip\noindent\textit{Gradual: Ps. 78:9 (Vg.; Heb. 79), Ps. 78:10 (Vg.; Heb. 79)}\par\nobreak
\begin{itemize}[leftmargin=1.2em,itemsep=1pt,topsep=1pt,parsep=0pt]
\item \textbf{Historical setting}. As under Alleluia: Ps. 9:5, 10 (Vg./Heb.).
\item \textbf{Composition}. As under Alleluia: Ps. 9:5, 10 (Vg./Heb.).
\end{itemize}
\par\smallskip\noindent\textit{Gospel: Luke 5:1--11}\par\nobreak
\begin{itemize}[leftmargin=1.2em,itemsep=1pt,topsep=1pt,parsep=0pt]
\item \textbf{Narrated event}. No date: the chronology corpus has not yet dated this (research pending).
\item \textbf{Composition} (disputed). The Gospel of St Luke: c.~A.D.~70 (``About the year 70''; The Catholic Encyclopedia, Volume 14 (New York, 1912)); Before the end of the Roman imprisonment, when the Acts was finished (De auctore, de tempore compositionis et de historica veritate Evangeliorum secundum Marcum et secundum Lucam (1912), responsa-i-ix).
\end{itemize}
\par\smallskip\noindent\textit{Epistle: Rom. 8:18--23}\par\nobreak
\begin{itemize}[leftmargin=1.2em,itemsep=1pt,topsep=1pt,parsep=0pt]
\item \textbf{Composition} (disputed). The Epistle to the Romans: A.D.~57 (``(Romans), 57''; The Catholic Encyclopedia, Volume 11 (New York, 1911)); A.D.~58 (``the Epistle to the Romans (58)''; The Catholic Encyclopedia, Volume 5 (New York, 1909)).
\end{itemize}
\par}
